% 学習結果説明（手順3a：用語集と図の見方）
% results フォルダで「lualatex 学習結果説明_手順3a.tex」を2回実行すると PDF を作り直せる（TeX Live の LuaLaTeX を使用）。
% 図は results/step3a/ の PNG を相対パスで読み込む。
\documentclass[paper=a4paper,fontsize=10.5pt,line_length=44zw,number_of_lines=38]{jlreq}
\usepackage[haranoaji]{luatexja-preset}
\usepackage{amsmath}
\usepackage{graphicx}
\usepackage{booktabs}
\usepackage{tabularx}
\usepackage{enumitem}
\usepackage{float}
\usepackage{xcolor}
\usepackage[hidelinks,bookmarksnumbered]{hyperref}
\graphicspath{{step3a/}}

\newcommand{\file}[1]{\texttt{\detokenize{#1}}}
\newcommand{\gd}{\dot{\gamma}}
\newcommand{\PI}{P3\_1}
\newcommand{\PII}{P3\_2}
\newcommand{\Tff}{T_{\mathrm{ff}}}
\newcommand{\Tset}{T_{\mathrm{set}}}
% 用語集：「・用語：説明文」
\newenvironment{termlist}{\begin{itemize}[label=・,leftmargin=1\zw,labelwidth=1\zw,labelsep=0pt,itemsep=0.4ex,topsep=0.6ex,parsep=0pt]}{\end{itemize}}
\newcommand{\term}[1]{\item \textbf{#1}：}
% 図の見方の項目
\newenvironment{points}{\begin{itemize}[label=$\triangleright$,leftmargin=1.5\zw,itemsep=0.3ex,topsep=0.5ex,parsep=0pt]}{\end{itemize}}
\newcommand{\pt}[1]{\item \textbf{#1}：}
\renewcommand{\arraystretch}{1.25}

\title{樹脂粘度モデル開発　手順3aの学習結果の説明\\[0.6ex]{\large 用語集とグラフの見方（機械学習・統計を初めて学ぶ人向け）}}
\author{}
\date{@@DATE@@}

\begin{document}
\maketitle
{\small\tableofcontents}

% =====================================================================
\section{この文書について}

この文書は、\file{results/step3a/} にある手順3aの学習結果（図と表）を、機械学習や統計を初めて学ぶ人でも読めるように説明したものです。射出成形と樹脂の粘度の基本は知っていることを前提にしています。手順2の説明書（\file{学習結果説明.pdf}）と同じ形で、重なる用語は短く繰り返しています。

\begin{points}
\pt{対象の結果}\file{results/step3a/}：新しいデータ（PP 再々実験と粉砕材混合 PP）で、流動中樹脂圧力の2列を説明変数に使えるかを調べた結果です（図6枚、CSV 11 個）。
\pt{読み方}2章で何をしたかをつかみ、分からない言葉は3章の用語集で調べてください。4章で図の見方、5章で表の見方、6章で結果のまとめと、質問への答えを説明します。
\pt{数値について}文中の数値は、この文書を作った時点の results の表から取り出したものです。
\pt{元のノート}\file{01_機械学習コードフォルダ} の \file{手順3a_新データ_圧力列の検討.ipynb} です。
\pt{記号表}記号・モデル名（A0、L1 など）・列名の一覧は、\file{学習結果説明_記号表.pdf} にあります（3a の用語集で説明済みの記号も含みます）。
\pt{作り直し方}\file{results} フォルダで \texttt{lualatex 学習結果説明\_手順3a.tex} を2回実行すると、この PDF を作り直せます。
\end{points}

% =====================================================================
\section{何をしたか（概要）}

\subsection{目的}
新しいデータには、流動中の樹脂の圧力が \PI\ と \PII\ の2列あります。\PI\ は「センサ4に樹脂が届いた時刻のセンサ3の圧力」、\PII\ は「センサ4の出力が最大になった時刻のセンサ3の圧力」です。この2列を、粘度を予測するモデルの入力に使ってよいかを、次の5つの観点で確かめました。

\begin{table}[H]
\caption{確かめた5つの観点}
\begin{tabularx}{\linewidth}{@{}>{\raggedright\arraybackslash}p{9\zw}>{\raggedright\arraybackslash}p{16\zw}X@{}}
\toprule
観点 & 確かめること & 見る場所 \\
\midrule
1 リーク & 圧力から粘度が計算で決まってしまわないか（線形回帰で $R^2 \approx 1$、係数が $\pm 1$ に近くならないか）。 & 表\ref{tab:ols}、図1 \\
2 予測のときに使えるか & 成形条件から粘度を予測するとき、その値を事前に知ることができるか。 & 6章 \\
3 検証データでの精度 & 低速側と高速側に分けた RMSE、95\% 予測区間の被覆率、NLPD。 & 図1・図2、表\ref{tab:sv_pp}・\ref{tab:sv_mix} \\
4 物理的な妥当さ & どの入力がよく効いているか（ARD）。学習範囲の外で予測の幅が広がるか。 & 図3・図4 \\
5 ノイズ分散 & 同じ条件のショット間のばらつき（純粋誤差）と比べて、モデルのノイズ分散はどうか。 & 表\ref{tab:noise} \\
\bottomrule
\end{tabularx}
\end{table}

あわせて、3b 以降の準備として、低速側への外挿のために GPR に線形の傾向を入れる方法を試算しました（図5）。

\subsection{データ}
新しいデータは、\file{02_機械学習データ/NR500/}（学習側）と \file{03_予測検証用テストデータ/NR500/}（検証側）にあります。1行が1ショットです。

\begin{points}
\pt{材料}PP と、PP に一次成形の PP（粉砕材）を 25 wt\% 加えた粉砕材混合 PP の2つです。
\pt{学習データ（5data）}PP は @@PP_NTR@@ 行（@@PP_NCTR@@ 条件）、粉砕材混合 PP は @@MIX_NTR@@ 行（@@MIX_NCTR@@ 条件）。各条件 4〜5 ショットです。
\pt{検証データ（3data）}PP は @@PP_NTE@@ 行（@@PP_NCTE@@ 条件）、粉砕材混合 PP は @@MIX_NTE@@ 行（@@MIX_NCTE@@ 条件）。各条件 1〜3 ショットです。学習データと条件は重なりません。検証データは、射出速度の両端にあります（図0）。低速側は 3・5 mm/s（厚み 1.5・2.0 mm のみ）、高速側は 280・320 mm/s（全厚み）です。
\pt{Alldata}学習と検証をまとめたファイルで、従来のランダム分割（行を2割ランダムにテストにする）で評価するために使いました。PP の Alldata は、各条件 10 ショット全部を含みます。
\pt{外れ値}「外れ値なし」を主に使い、「外れ値あり」は結果が変わらないかの確認だけに使いました。PP は、圧力の最大値が検出できなかった厚み 1.5 mm・射出速度 3 mm/s の3条件を除いたファイル（no1.5-3.0）を使いました。
\pt{圧力の単位}CSV の \PI・\PII\ は MPa です。学習では Pa に直してから $\ln$ をとりました（6章の質問への答えを参照）。
\end{points}

\begin{table}[H]
\caption{使ったデータの概要（外れ値なし。\file{data_summary.csv}）}
\begin{tabularx}{\linewidth}{@{}llrrrll@{}}
\toprule
材料 & 組 & 行数 & 条件数 & ショット/条件 & $\gd$ [1/s] & $\eta$ [Pa$\cdot$s] \\
\midrule
@@TABLE_DATA@@
\bottomrule
\end{tabularx}
\end{table}

\subsection{比べたモデル}
目的変数（予測したい量）は、どのモデルも粘度の自然対数 $\ln\eta$ です。入力はすべて $\ln$ をとった値です（$\Tff$ は K で）。

\begin{table}[H]
\caption{比べた7つのモデル}
\begin{tabularx}{\linewidth}{@{}l>{\raggedright\arraybackslash}p{14\zw}X@{}}
\toprule
記号 & モデルと入力 & 意味 \\
\midrule
L0 & 線形回帰\newline $\gd$，$H$ & 圧力を使わない基準。 \\
L1 & 線形回帰\newline \PI，$\gd$，$H$ & 粘度の算出式に近い形。リークの確認用。 \\
L2 & 線形回帰\newline \PII，$\gd$，$H$ & \PII\ がどれだけ役に立つかの確認用。 \\
A0 & GPR\newline $\gd$，$H$，$\Tff$ & 圧力を使わない GPR。 \\
A1 & GPR\newline $\gd$，$H$，$\Tff$，\PI & A0 に \PI\ を足したもの。 \\
A2 & GPR\newline $\gd$，$H$，$\Tff$，\PII & A0 に \PII\ を足したもの。 \\
A1$'$ & GPR\newline $\gd$，$\Tff$，\PI & 現行ノートと同じ形（厚み $H$ がない）。 \\
\bottomrule
\end{tabularx}
\end{table}

GPR はすべて、ARD 付き Matern 5/2 カーネルにノイズ項を足したものです。

\subsection{データの分け方}
\begin{table}[H]
\caption{学習データとテストデータの分け方}
\begin{tabularx}{\linewidth}{@{}lX@{}}
\toprule
記号 & 分け方 \\
\midrule
S\_V & 5data で学習し、3data で検証する。射出速度の両端の外挿にあたる。低速側（3・5 mm/s）と高速側（280・320 mm/s）に分けて評価する。条件は学習と検証で重ならない。 \\
S\_R1 & Alldata の行を、ランダムに2割テストにする（従来と同じ分け方）。同じ条件のほかのショットが学習にも入る。 \\
S\_R2 & 5data と 3data をまとめた行を、ランダムに2割テストにする（PP のみ。S\_R1 とのデータ量の違いを見る）。 \\
LOCO & 条件を1つずつ抜いて、その条件を予測する（条件ごとに抜く交差検証）。 \\
\bottomrule
\end{tabularx}
\end{table}

\subsection{そのほかの設定}
\begin{points}
\pt{標準化と中心化}学習データだけで計算しました（旧ノートは全行で計算していました）。
\pt{ハイパーパラメータ}最尤推定で、1組だけ求めました。ノイズ分散の下限は $10^{-6}$ に下げました（旧ノートは $10^{-3}$ で、下限に張り付いていました）。
\pt{乱数}シードは 42 に固定しました。
\end{points}

% =====================================================================
\section{用語集}

「・用語：説明」の形で、分野ごとに並べました。

\subsection{粘度と測定に関する用語}
\begin{termlist}
\term{見かけの粘度 $\eta$ [Pa$\cdot$s]}樹脂の流れにくさ。キャビティ内の2点の圧力センサの値から計算した値で、せん断発熱や冷却の影響も含んだ「見かけ」の値です。
\term{見かけのせん断速度 $\gd$ [1/s]}流れの中で、樹脂がずれるように変形する速さ。$\gd = 6L/(H\,\Delta t)$ で、センサ2点の間を樹脂が通る時間 $\Delta t$ から計算します。ショットごとに測った値です。
\term{圧力 \PI・\PII}どちらも、センサ3の圧力です。\PI\ は、センサ4に樹脂が届いた時刻の値で、粘度の計算に使われています。\PII\ は、センサ4の出力が最大になった時刻の値で、\PI\ の約 4〜25 倍あり、粘度の計算には使われていません。
\term{MPa と Pa}圧力の単位。1 MPa = $10^6$ Pa。CSV の圧力は MPa です。
\term{射出速度 [mm/s]}成形機に設定した射出の速さ。成形条件として事前に決める値です（射出率 [cm$^3$/s] はこの 0.531 倍）。
\term{加熱シリンダ設定温度 $\Tset$}成形機に設定した樹脂の温度。190・210・230 ℃ の3水準です。
\term{流動先端温度 $\Tff$}センサの間を流れている樹脂の先端の温度（実測値）。流れの速さによって上下するので、流れの「結果」でもあります。
\term{キャビティ厚み $H$ [mm]}平板のキャビティの厚み。0.5・1.0・1.5・2.0 mm の4水準です。
\term{条件とショット}条件は、設定温度・厚み・射出速度の組み合わせです。同じ条件で複数回成形しており、1回の成形をショットと呼びます。
\term{$\ln$}自然対数（$e \approx 2.718$ を底とする対数）。粘度は桁で変わるので、計算は対数で行います。$\ln\eta$ の差 0.1 は、粘度の約 1.1 倍の違いに当たります。
\end{termlist}

\subsection{データの扱いに関する用語}
\begin{termlist}
\term{目的変数と説明変数}目的変数は予測したい量（ここでは $\ln\eta$）、説明変数（入力）は予測の手がかりにする量です。
\term{学習データとテストデータ（検証データ）}学習データは、モデルの中の値を決めるときに使うデータ。テストデータ（検証データ）は、学習には使わず、できたモデルの予測がどれだけ当たるかを確かめるためだけに使うデータです。
\term{5data・3data・Alldata}5data は学習用（各条件 4〜5 ショット）、3data は検証用（各条件 1〜3 ショット）のファイルです。Alldata は、これらをまとめたファイルです（PP の Alldata は、各条件 10 ショット全部）。
\term{外れ値}ほかの値から大きく外れた値。測定の失敗などが考えられます。「外れ値なし」のファイルは、あらかじめ決めた基準（中央値から 20\% 以上離れた値など）で外れ値を除いたものです。
\term{ランダム分割と条件単位の分割}ランダム分割は、行をランダムに学習とテストに分けます。同じ条件のほかのショットが学習にも入るので、「ほとんど同じデータでの答え合わせ」になり、評価が甘くなります。条件単位の分割（S\_V など）は、条件ごとまとめて分けるので、評価が厳しくなります。
\term{内挿と外挿}学習データの範囲の内側を予測するのが内挿、外側を予測するのが外挿です。外挿のほうが難しくなります。S\_V の検証は外挿にあたります。
\term{標準化}入力の値から平均を引き、標準偏差で割って、平均 0・ばらつき 1 にそろえること。単位や大きさの違う入力を公平に扱うためです。
\term{中心化}値から平均を引いて、平均を 0 にそろえること。
\term{リーク（目的変数のリーク）}答え（目的変数）を計算するのに使った値が入力に紛れ込み、モデルが答えを「計算」できてしまうこと。精度は高く見えますが、実際に予測したい場面ではその入力を先に知ることができないので、役に立ちません。\PI\ がこれに当たります。
\end{termlist}

\subsection{モデルに関する用語}
\begin{termlist}
\term{線形回帰}入力に重みを掛けて足し合わせ、定数（切片）を加えた式で予測する、最も単純な回帰。重みは最小二乗法（予測と実測の差の2乗の合計が最小になるように決める方法）で決めます。$\ln$ をとった入力に対する線形回帰の係数は、「入力が1\% 増えると、粘度が何\% 変わるか」を表します。
\term{ガウス過程回帰（GPR）}「似た入力には似た出力が対応する」という考えで、データの間を滑らかにつなぐ回帰の方法。予測値と一緒に、予測のばらつき（不確かさ）も出せます。近くにデータが少ない所では、ばらつきが大きくなります。
\term{カーネル}2つの入力がどれくらい「似ているか」を決める関数。GPR の性質はカーネルで決まります。Matern 5/2 カーネルは、入力が近いほど似ているとみなす、ほどよく滑らかなカーネルです。
\term{ARD と長さスケール}入力ごとに別々の長さスケールを持たせる仕組みが ARD。長さスケールは、入力がどれくらい変わると出力が大きく変わるかの目安（標準化した入力に対する値）で、小さいほどその入力がよく効いています。この文書では、その逆数（$=1/$長さスケール）を「感度」として図に描いています。長さスケールが探索範囲の上限（100）に張り付いているときは、その入力がほとんど効いていないことを表します。
\term{ノイズ分散と潜在関数}GPR は、観測値を「滑らかな関数（潜在関数）＋ノイズ」とみなします。ノイズ分散は、同じ条件でもショットごとにばらつく大きさ（成形や測定のばらつき）です。予測の標準偏差には、潜在関数の不確かさとノイズの両方が入っています。
\term{定数平均}GPR の標準的な設定で、「データがない所では、学習データの平均値に戻る」という性質を持ちます。学習範囲の外で外れる原因になります。
\term{線形カーネル（DotProduct）}カーネルに足すと、GPR が「直線的な傾向」も表せるようになります。傾きもデータから同時に決まり、傾きの不確かさも予測の幅に入ります。
\term{残差}実測と予測の差。「線形回帰の残差を GPR で予測する」は、線形回帰で説明しきれなかった残りだけを GPR に任せる方法です。
\term{ハイパーパラメータ}モデルの形を決める設定値。GPR では、振幅・長さスケール・ノイズ分散です。データから推定します（最尤推定：データが最ももっともらしくなる1組を選ぶ方法）。
\end{termlist}

\subsection{評価の指標}
\begin{termlist}
\term{決定係数 $R^2$}予測が、データのばらつきをどれだけ説明できたかを表す指標。1 なら完全に一致、0 なら「いつも平均値を答える」のと同じ、マイナスならそれより悪いことを表します。低速側のようにテストデータの値の範囲が狭いと、同じ誤差でも小さくなり、大きなマイナスにもなります。
\term{RMSE（二乗平均平方根誤差）}予測と実測の差の、平均的な大きさ。$\ln\eta$ の RMSE が 0.05 なら、粘度で約 5\% の誤差に当たります。小さいほどよい指標です。
\term{予測区間と被覆率}予測区間は、新しく測った値が 95\% の確率で入ると予測した範囲。被覆率は、テストデータのうち、予測区間に実測値が入った割合です。0.95 に近いのが理想で、小さすぎると区間が狭すぎ（自信過剰）、1.00 は区間が広すぎる可能性があります。
\term{NLPD（負の対数予測密度）}予測を「平均と標準偏差の正規分布」とみたとき、実測値がどれくらい「ありそう」だったかを表す指標。当たり具合と、予測のばらつきの見積もりの適切さを、まとめて評価します。小さい（マイナスに大きい）ほどよい指標です。区間が狭すぎて外れると大きくなります。
\term{純粋誤差}同じ条件のショット間のばらつき（標準偏差）。モデルでは減らせないばらつきの目安です。ここでは $\ln\eta$ のばらつきで、「外れ値なし」のデータは、中央値から 20\% 以上離れた値を除いているので、小さめに出ます。
\term{LOCO（条件を1つずつ抜く交差検証）}条件を1つ抜いて学習し、抜いた条件を予測することを、全条件で繰り返す評価。ここでは、ハイパーパラメータを全データで決めて固定し、閉じた式で計算しました（近似）。抜くのは1条件だけで、同じ厚み・射出速度のほかの温度は学習に残るので、速度の外挿には厳しくない評価です。
\term{パリティ図}横軸に実測値、縦軸に予測値をとって点を打った図。点が対角線（理想線）に乗るほど、予測が当たっています。縦棒は 95\% 予測区間です。
\term{対数軸}目盛りが 1，10，100 のように、桁ごとに等間隔になる軸。桁で変わる量（粘度やせん断速度）を見やすくします。
\end{termlist}

% =====================================================================
\section{図の見方}

図ごとに、「何を示す図か」「軸と記号」「見るポイント」「今回の結果」の順に説明します。図のファイルは \file{results/step3a/} にあります。図の番号は、ファイル名（\file{fig0_…}〜\file{fig5_…}）に合わせて、0 から数えています。
\setcounter{figure}{-1}

% ---------------------------------------------------------------------
\subsection{図0：データの全体像}
\begin{figure}[H]
\centering
\includegraphics[width=\linewidth]{fig0_data_overview.png}
\caption{学習データ（塗りつぶし）と検証データ（白抜き）の位置}
\end{figure}
\begin{points}
\pt{ファイル}\file{results/step3a/fig0_data_overview.png}
\pt{何を示す図か}学習データと検証データが、せん断速度と粘度の上でどこにあるかを示す、データそのものの図です。上段が PP、下段が粉砕材混合 PP、列が厚みです。
\pt{軸}横軸は見かけのせん断速度 $\gd$ [1/s]、縦軸は見かけの粘度 $\eta$ [Pa$\cdot$s]（どちらも対数軸）です。
\pt{記号}塗りつぶしの丸が学習データ、白抜きの四角が検証データ。色は設定温度（青 190 ℃、緑 210 ℃、赤 230 ℃）です。灰色の帯は、学習データのせん断速度の範囲です。
\pt{見るポイント}白抜きの四角が、灰色の帯の外（左が低速側、右が高速側）にあるか。
\pt{今回の結果}検証データは、学習データの範囲の外にあります。低速側は、せん断速度が学習の約 @@GD_LOW_PP@@ 分の 1 まで小さく（学習 @@PP_GD_TR@@ に対し検証 @@PP_GD_TE@@ 1/s）、粘度は学習の最大の約 @@ETA_RATIO_PP@@ 倍（学習 @@PP_ETA_TR@@ に対し検証 @@PP_ETA_TE@@ Pa$\cdot$s）に達します。高速側は、学習範囲を少し超える程度です。
\end{points}

% ---------------------------------------------------------------------
\subsection{図1：パリティ図（5data → 3data）}
\begin{figure}[H]
\centering
\includegraphics[width=\linewidth]{fig1_parity_S_V.png}
\caption{検証データでのパリティ図（上段 PP、下段 粉砕材混合 PP。列は L1，L2，A0，A1，A2，A1$'$）}
\end{figure}
\begin{points}
\pt{ファイル}\file{results/step3a/fig1_parity_S_V.png}
\pt{何を示す図か}6つのモデル（列）の予測が、検証データでどれだけ当たったかを、2つの材料（行）で比べる図です。
\pt{軸}横軸は実測の $\ln\eta$、縦軸は予測の $\ln\eta$ です。
\pt{記号}点が1ショット、青が低速側、赤が高速側、灰色の縦線が 95\% 予測区間、黒い破線が理想線です。各パネルの上に、モデルと、低速側・高速側の RMSE を示しました。
\pt{見るポイント}点が破線に沿っているか。灰色の縦線が破線にかかっているか（被覆率）。
\pt{今回の結果（リーク）}\PI\ を入れた L1 と A1 は、低速側も高速側も破線にほぼ重なります（RMSE は L1 で @@L1_PP_LOW@@ と @@L1_PP_HIGH@@（PP）、A1 で @@A1_PP_LOW@@ と @@A1_PP_HIGH@@）。これは、\PI\ が粘度の算出式そのものだからで、モデルが算出式を覚えているだけです。
\pt{今回の結果（\PII）}L2 は、\PII\ を入れても、圧力を使わない L0 より良くなりません（低速側の RMSE は PP で L2 が @@L2_PP_LOW@@、L0 が @@L0_PP_LOW@@）。GPR の A2 も、圧力を使わない A0 とほとんど同じです（低速側 PP で @@A2_PP_LOW@@ と @@A0_PP_LOW@@。粉砕材混合 PP では @@A2_MIX_LOW@@ と @@A0_MIX_LOW@@ で、わずかに良い程度）。
\pt{今回の結果（低速側）}A0 と A2 は、高速側（赤）は破線に乗りますが、低速側（青）は破線から大きく下にずれます（PP の RMSE：高速側 @@A0_PP_HIGH@@、低速側 @@A0_PP_LOW@@）。学習データより粘度が大きい領域で、GPR が学習データの平均に引き戻され、粘度を小さく見積もってしまうためです。縦線は長いのですが、破線まで届いていません（被覆率は低速側で @@A0_PP_COV_LOW@@）。
\pt{今回の結果（A1$'$）}厚み $H$ を入れていない A1$'$ は、高速側は当たりますが、低速側は外れます。圧力を外すと厚みの情報がなくなるので、$H$ は入力に必要です。
\pt{今回の結果（PP の高速側の NLPD）}PP の高速側の NLPD は、A0 で @@A0_PP_NLPD_HIGH@@ と悪く見えますが、これはほぼ、190 ℃・厚み 0.5 mm・射出速度 320 mm/s の条件の @@ANOM_N@@ ショットで決まっています（表\ref{tab:anom_cond}）。この条件は、同じ速度の 210・230 ℃ に比べて、粘度が約 @@ANOM_RATIO@@ 倍と外れています。予測の標準偏差が小さい所で大きく外れると、NLPD が急に大きくなるためです。この条件を除くと、PP の高速側の NLPD は @@A0_PP_NLPD_HIGH_X@@、RMSE は @@A0_PP_RMSE_HIGH_X@@ になります（表\ref{tab:anom}）。この条件の扱いは、3b の感度確認で確かめます。
\end{points}

\begin{table}[H]
\caption{線形回帰の係数と検証の $R^2$（5data で当てはめ。\file{ols_coefficients.csv}、\file{metrics.csv}）}\label{tab:ols}
\begin{tabularx}{\linewidth}{@{}l>{\raggedright\arraybackslash}p{6\zw}rrrrr@{}}
\toprule
材料 & モデル & 切片 & $\ln P$ の係数 & $\ln\gd$ の係数 & $\ln H$ の係数 & 検証 $R^2$ \\
\midrule
@@TABLE_OLS@@
\bottomrule
\end{tabularx}
\end{table}

\begin{table}[H]
\caption{検証データ（5data $\to$ 3data）での指標：PP（$\ln\eta$。\file{metrics.csv}）}\label{tab:sv_pp}
\begin{tabularx}{\linewidth}{@{}X*{6}{r}@{}}
\toprule
 & \multicolumn{2}{c}{RMSE} & \multicolumn{2}{c}{被覆率} & \multicolumn{2}{c}{NLPD} \\
\cmidrule(lr){2-3}\cmidrule(lr){4-5}\cmidrule(l){6-7}
モデル & 低速側 & 高速側 & 低速側 & 高速側 & 低速側 & 高速側 \\
\midrule
@@TABLE_SV_PP@@
\bottomrule
\end{tabularx}
\end{table}

\begin{table}[H]
\caption{検証データ（5data $\to$ 3data）での指標：粉砕材混合 PP}\label{tab:sv_mix}
\begin{tabularx}{\linewidth}{@{}X*{6}{r}@{}}
\toprule
 & \multicolumn{2}{c}{RMSE} & \multicolumn{2}{c}{被覆率} & \multicolumn{2}{c}{NLPD} \\
\cmidrule(lr){2-3}\cmidrule(lr){4-5}\cmidrule(l){6-7}
モデル & 低速側 & 高速側 & 低速側 & 高速側 & 低速側 & 高速側 \\
\midrule
@@TABLE_SV_MIX@@
\bottomrule
\end{tabularx}
\end{table}

\begin{table}[H]
\caption{検証データの厚み 0.5 mm・射出速度 320 mm/s（外れ値なし。\file{anomaly_condition.csv}）}\label{tab:anom_cond}
\begin{tabularx}{\linewidth}{@{}lrrrrr@{}}
\toprule
材料 & 設定温度 [℃] & ショット数 & $\gd$ の平均 [1/s] & $\eta$ の平均 [Pa$\cdot$s] & $\Tff$ の平均 [℃] \\
\midrule
@@TABLE_ANOM_COND@@
\bottomrule
\end{tabularx}
\end{table}

\begin{table}[H]
\caption{高速側の指標：190 ℃・0.5 mm・320 mm/s の条件を除いた場合（A0。$\ln\eta$。\file{anomaly_check.csv}）}\label{tab:anom}
\begin{tabularx}{\linewidth}{@{}l>{\raggedright\arraybackslash}p{12\zw}rrrr@{}}
\toprule
材料 & 範囲 & ショット数 & RMSE & 被覆率 & NLPD \\
\midrule
@@TABLE_ANOM@@
\bottomrule
\end{tabularx}
\end{table}

% ---------------------------------------------------------------------
\subsection{図2：分け方による評価の違い}
\begin{figure}[H]
\centering
\includegraphics[width=0.92\linewidth]{fig2_metrics_by_split.png}
\caption{分け方ごとの RMSE（上段）と NLPD（下段）}
\end{figure}
\begin{points}
\pt{ファイル}\file{results/step3a/fig2_metrics_by_split.png}
\pt{何を示す図か}同じモデルでも、データの分け方で評価がどれだけ変わるかを比べる図です。左が PP、右が粉砕材混合 PP。
\pt{軸}横軸はモデル、縦軸は上段が RMSE（$\ln\eta$。対数軸）、下段が NLPD（小さいほど良い）です。
\pt{記号}棒の色が分け方です。青：S\_V の低速側、赤：S\_V の高速側、緑：ランダム 20\%（Alldata）、黄緑：ランダム 20\%（5data＋3data。PP のみ）。
\pt{見るポイント}同じモデルで、緑・黄緑の棒が、青・赤の棒より大幅に低くないか。
\pt{今回の結果}ランダム 20\% の RMSE は、A0 で PP @@A0_PP_R1@@、粉砕材混合 PP @@A0_MIX_R1@@ と、とても小さい値です。PP では、S\_V の高速側の約 1/@@RATIO_HIGH@@、低速側の約 1/@@RATIO_LOW@@ にすぎません。ランダム分割では、テストの条件のほぼすべてが学習にも入っているからです（PP の Alldata：@@OVERLAP_PP@@ 条件、粉砕材混合 PP：@@OVERLAP_MIX@@ 条件）。従来の評価は、モデルの本当の力を、かなり甘く見せていました。
\end{points}

\begin{table}[H]
\caption{分け方ごとの RMSE（$\ln\eta$。外れ値なし）}
\begin{tabularx}{\linewidth}{@{}X*{7}{r}@{}}
\toprule
 & \multicolumn{4}{c}{PP} & \multicolumn{3}{c}{粉砕材混合 PP} \\
\cmidrule(lr){2-5}\cmidrule(l){6-8}
モデル & S\_V 低速 & S\_V 高速 & S\_R1 & S\_R2 & S\_V 低速 & S\_V 高速 & S\_R1 \\
\midrule
@@TABLE_RANDOM@@
\bottomrule
\end{tabularx}
\end{table}

\begin{table}[H]
\caption{条件を1つずつ抜く交差検証（LOCO）の結果（$\ln\eta$。\file{loco_metrics.csv}）}
\begin{tabularx}{\linewidth}{@{}X*{4}{r}@{}}
\toprule
 & \multicolumn{2}{c}{PP} & \multicolumn{2}{c}{粉砕材混合 PP} \\
\cmidrule(lr){2-3}\cmidrule(l){4-5}
モデル & RMSE & 被覆率 & RMSE & 被覆率 \\
\midrule
@@TABLE_LOCO@@
\bottomrule
\end{tabularx}
\end{table}

LOCO では、A0 でも RMSE が小さく出ます。抜くのは1条件だけで、同じ厚み・射出速度のほかの温度の条件が学習に残るためで、低速側の外挿の厳しさは測れません。

% ---------------------------------------------------------------------
\subsection{図3：ARD の感度}
\begin{figure}[H]
\centering
\includegraphics[width=\linewidth]{fig3_ard_sensitivity.png}
\caption{GPR の入力ごとの感度（1／長さスケール）}
\end{figure}
\begin{points}
\pt{ファイル}\file{results/step3a/fig3_ard_sensitivity.png}
\pt{何を示す図か}GPR が、どの入力をよく使っているかを示す図です。左が PP、右が粉砕材混合 PP。5data で学習したモデルの値です。
\pt{軸}横軸は入力（せん断速度 $\gd$、厚み $H$、流動先端温度 $\Tff$、\PI、\PII）、縦軸は感度（長さスケールの逆数。標準化した入力に対する値）です。
\pt{記号}棒の色がモデル（A0，A1，A2，A1$'$）です。棒がない入力は、そのモデルに入っていません。
\pt{見るポイント}棒が高い入力ほど、よく効いています。棒がほぼ 0 の入力は、効いていません。
\pt{今回の結果}A0 では、$\gd$ の長さスケールが @@A0_PP_LSG@@（PP）・@@A0_MIX_LSG@@（粉砕材混合 PP）と小さく、圧倒的に効いています。$H$ は @@A0_PP_LSH@@・@@A0_MIX_LSH@@、$\Tff$ は @@A0_PP_LST@@・@@A0_MIX_LST@@ と大きく、ほとんど効いていません。A2 の \PII\ の長さスケールは、PP で @@A2_PP_LSP@@（上限 100 に張り付き）、粉砕材混合 PP で @@A2_MIX_LSP@@ で、ほとんど使われていません。A1 は、どの入力も長さスケールが大きく（約 @@A1_LS_MIN@@ 以上）、ほぼ直線的な関数、つまり算出式そのものを表しています。
\end{points}

% ---------------------------------------------------------------------
\subsection{図4：予測の幅と射出速度}
\begin{figure}[H]
\centering
\includegraphics[width=\linewidth]{fig4_predictive_sd_vs_speed.png}
\caption{検証データの予測標準偏差（ノイズ込み）の平均と、射出速度}
\end{figure}
\begin{points}
\pt{ファイル}\file{results/step3a/fig4_predictive_sd_vs_speed.png}
\pt{何を示す図か}学習範囲の外で、GPR が「分からない」と予測の幅を広げるかを確かめる図です。
\pt{軸}横軸は射出速度 [mm/s]（対数軸）、縦軸は検証データの予測標準偏差（$\ln\eta$）の平均（対数軸）です。
\pt{記号}線の色がモデル、灰色の帯が学習データの射出速度の範囲（10〜240 mm/s）です。
\pt{見るポイント}帯の外（特に左端の低速側）で、線が上がっているか。
\pt{今回の結果}圧力を使わないモデル（A0）の予測標準偏差は、低速側ほど大きくなります（低速側の平均は PP で @@A0_PP_SD_LOW@@、高速側は @@A0_PP_SD_HIGH@@）。学習範囲の外で幅を広げる性質は、働いています。ただし、実際の誤差（RMSE @@A0_PP_LOW@@）より小さく、足りません（図1）。A1 は、算出式を覚えているので、どの速度でも幅がとても小さくなります。
\end{points}

\begin{table}[H]
\caption{ノイズ分散と純粋誤差（5data で学習。$\ln\eta$ の標準偏差。\file{kernels.csv}）}\label{tab:noise}
\begin{tabularx}{\linewidth}{@{}l>{\raggedright\arraybackslash}p{4\zw}rrrl@{}}
\toprule
材料 & モデル & ノイズの標準偏差 & 純粋誤差 & $\ln\gd$ で説明した残り & 下限に張り付き \\
\midrule
@@TABLE_NOISE@@
\bottomrule
\end{tabularx}
\end{table}

ノイズの標準偏差は、同じ条件の中のばらつき（純粋誤差）よりずっと小さく、「$\ln\gd$ で説明した残り」とほぼ同じです。同じ条件の中では、ショットごとの $\ln\eta$ と $\ln\gd$ が強く連動していて（相関 約 $-0.98$）、ショットごとの $\gd$ を入力にすると、ばらつきの大部分が説明されるためです。成形条件（射出速度など、事前に決める値）から見たばらつきは、これより大きくなります。3b で調べます。

% ---------------------------------------------------------------------
\subsection{図5：線形の傾向を入れた試算}
\begin{figure}[H]
\centering
\includegraphics[width=0.9\linewidth]{fig5_linear_trend_trial.png}
\caption{低速側への外挿の試算（上段 PP、下段 粉砕材混合 PP）}
\end{figure}
\begin{points}
\pt{ファイル}\file{results/step3a/fig5_linear_trend_trial.png}
\pt{何を示す図か}低速側で外れる問題への対策として、GPR に線形の傾向を入れる2つの方法を、図1と同じ検証データで試した図です。列は、定数平均の GPR（A0。比較用）、(a) 線形回帰の残差を GPR で予測、(b) カーネルに線形項を足す、です。
\pt{軸と記号}図1と同じです（青が低速側、赤が高速側）。
\pt{見るポイント}青い点が破線に近づいたか。縦線が破線にかかっているか。
\pt{今回の結果}低速側の RMSE は、PP で A0 の @@TR_PP_A0_LOW@@ から (a) @@TR_PP_A_LOW@@、(b) @@TR_PP_B_LOW@@ へ、粉砕材混合 PP で @@TR_MIX_A0_LOW@@ から @@TR_MIX_A_LOW@@、@@TR_MIX_B_LOW@@ へ下がりました。低速側の被覆率は (b) で PP @@TR_PP_B_COV_LOW@@、粉砕材混合 PP @@TR_MIX_B_COV_LOW@@ です。高速側は変わりません。3b 以降は、(b) を使います。
\end{points}

\begin{table}[H]
\caption{線形の傾向の試算（検証データ。\file{linear_trend_trial.csv}）}
\begin{tabularx}{\linewidth}{@{}l>{\raggedright\arraybackslash}p{10\zw}*{6}{r}@{}}
\toprule
 &  & \multicolumn{2}{c}{RMSE} & \multicolumn{2}{c}{被覆率} & \multicolumn{2}{c}{NLPD} \\
\cmidrule(lr){3-4}\cmidrule(lr){5-6}\cmidrule(l){7-8}
材料 & 方法 & 低速 & 高速 & 低速 & 高速 & 低速 & 高速 \\
\midrule
@@TABLE_TREND@@
\bottomrule
\end{tabularx}
\end{table}

% =====================================================================
\section{表（CSV）の見方}

\file{results/step3a/} の CSV ファイルの中身と、主な列の意味です。$\ln$ は自然対数です。

\begin{table}[H]
\caption{\file{results/step3a/} の表}
\begin{tabularx}{\linewidth}{@{}>{\raggedright\arraybackslash}p{13\zw}X@{}}
\toprule
ファイル & 中身と主な列 \\
\midrule
\file{metrics.csv} & 分け方・データ・モデル・評価の範囲ごとの指標。\file{split}（S\_V，S\_R1，S\_R2）、\file{data}（外れ値の なし・あり）、\file{subset}（全体・低速側・高速側）、\file{RMSE}、\file{R2}、\file{coverage95}（被覆率）、\file{NLPD}、\file{mean_sd}（予測標準偏差の平均）、\file{mean_sd_latent}（潜在関数の標準偏差の平均）、\file{noise_sd}。 \\
\file{kernels.csv} & GPR のハイパーパラメータ。\file{amplitude}（振幅）、\file{ls_*}（長さスケール）、\file{noise_sd}（ノイズの標準偏差）、\file{pure_sd_ln_eta}（純粋誤差）、\file{pure_sd_after_gd}（$\ln\gd$ で説明した残り）、\file{LML}（対数周辺尤度）、\file{noise_at_lower_bound}・\file{ls_at_bound}（下限・上限への張り付き）。 \\
\file{ols_coefficients.csv} & 線形回帰（L0，L1，L2）の係数。\file{intercept}（切片）、\file{ln_*}（各入力の係数）、\file{resid_sd}（残差の標準偏差）。 \\
\file{loco_metrics.csv} & 条件を1つずつ抜く交差検証の指標。\file{subset} は 全体・低速側・内側・高速側。 \\
\file{data_summary.csv} & 使ったファイルの概要（行数、条件数、ショット数、値の範囲）。 \\
\file{formula_check.csv} & 粘度の算出式の確認。\file{L推定 [mm]}（センサ間距離）、\file{c推定 [MPa]}、最大相対誤差、$\Delta t$ が 5 $\mu$s の倍数からずれる最大値。 \\
\file{pure_error.csv} & 同じ条件のショット間のばらつき。\file{sd_ln_eta}、\file{corr_with_ln_gd}（$\ln\gd$ との相関）、\file{var_after_gd}（$\ln\gd$ で説明した残りの分散）。 \\
\file{unit_check.csv} & 圧力の単位（MPa と Pa）で結果が変わらないことの確認。\file{予測平均の差の最大}、\file{切片の差（MPa − Pa）}。 \\
\file{linear_trend_trial.csv} & 線形の傾向を入れた試算の指標と、学習したカーネル。 \\
\file{anomaly_condition.csv} & 検証データの厚み 0.5 mm・射出速度 320 mm/s の、設定温度ごとの平均（ショット数、$\gd$、$\eta$、$\Tff$）。 \\
\file{anomaly_check.csv} & 190 ℃・0.5 mm・320 mm/s の条件を含めた場合・除いた場合の、高速側の指標（全モデル）。 \\
\bottomrule
\end{tabularx}
\end{table}

% =====================================================================
\section{結果のまとめ}

\subsection{観点ごとの結論}
\begin{enumerate}[itemsep=0.6ex,topsep=0.6ex]
\item \textbf{観点1（リーク）}：\PI\ はリークです。線形回帰（L1）で検証 $R^2$ が @@L1_PP_R2@@、係数が $\ln P$ で @@L1_PP_BP@@、$\ln\gd$ で @@L1_PP_BG@@、$\ln H$ で @@L1_PP_BH@@ と、算出式（$\pm 1$）に近い値になりました（表\ref{tab:ols}）。GPR に入れる（A1）と、誤差がほぼ 0 になります。
\item \textbf{\PII}：リークではありませんが、役に立ちません。L2 は L0 より良くならず、GPR（A2）は A0 とほぼ同じ精度でした。\PII\ の長さスケールは、PP では上限に張り付き、粉砕材混合 PP でも大きく、ほとんど使われていません。
\item \textbf{観点2（予測のときに使えるか）}：射出速度・厚み・設定温度は、成形条件として事前に決まります。\PI・\PII・$\Tff$・ショットごとの $\gd$ は、流動の結果として決まる値で、事前には分かりません。圧力は、使えたとしても、予測の場面では使いにくい入力です。
\item \textbf{観点3（精度）}：圧力を使わないモデル（A0）は、高速側（RMSE @@A0_PP_HIGH@@）は良いのですが、低速側（@@A0_PP_LOW@@）で大きく外れます。学習データより粘度が大きい領域への外挿で、定数平均の GPR が平均値に戻るためです。線形の傾向を入れると、低速側の RMSE が @@TR_PP_B_LOW@@ まで下がりました（図5）。なお、PP の高速側の NLPD は、1つの条件（190 ℃・0.5 mm・320 mm/s の @@ANOM_N@@ ショット）に大きく左右されます（表\ref{tab:anom}）。
\item \textbf{観点4（物理的な妥当さ）}：$\gd$ が圧倒的に効き、$\Tff$ はほとんど効きません（図3）。学習範囲の外で予測の幅は広がりますが、実際の誤差には足りません（図4）。設定温度ごとの予測曲線は、3b で調べます。
\item \textbf{観点5（ノイズ分散）}：ノイズの標準偏差（約 @@A0_PP_NSD@@）は、純粋誤差（約 @@PURE_PP@@）よりずっと小さく、ショットごとの $\gd$ を入力にしたことで、ばらつきの大部分が説明されたためです。
\item \textbf{評価の甘さ}：ランダム 20\% の RMSE は、PP の A0 で、S\_V の高速側の約 1/@@RATIO_HIGH@@、低速側の約 1/@@RATIO_LOW@@ です（図2）。評価は、条件単位の分け方で行う必要があります。
\item \textbf{結論}：圧力を使わない方針は妥当です。\PII\ は、使っても効果が小さく、予測の場面でも使いにくいので、参考として並べます。低速側への外挿のために、線形の傾向を入れた GPR を3b で使います。
\end{enumerate}

\subsection{質問への答え}
\begin{points}
\pt{圧力の単位（MPa と Pa）}CSV の圧力は MPa です。Pa に直してから $\ln$ をとりました。MPa のまま $\ln$ をとっても、GPR の予測は変わりません。$\ln P$ が全行で一律に $\ln 10^6 \approx 13.8$ だけずれるだけで、標準化（平均を引く）で消えるためです。確かめた結果、予測平均の差は最大でも @@UNIT_MAXDIFF@@ でした。線形回帰でも、予測は同じで、切片だけが（$\ln P$ の係数）$\times \ln 10^6$ だけ変わります（表\ref{tab:unit}）。
\pt{「線形回帰の残差を GPR で予測」は二重学習にならないか}ならない、と考えています。線形回帰は、学習データから傾きを決めます。GPR は、その残り（残差）だけを学びます。線形回帰で説明できた部分を、GPR が重ねて学ぶことはありません。予測は「線形回帰の予測 $+$ GPR の予測」の足し算で、テストデータは学習に使いません。ただし弱点が2つあります。線形の傾きは固定され、傾きの不確かさが予測区間に入らないこと。傾きと残りを別々に決めるので、全体として最適とは限らないことです。そこで、傾きも GPR の中で同時に決まり、傾きの不確かさも予測の幅に入る (b)（カーネルに線形項を足す）を、3b 以降に使います。
\end{points}

\begin{table}[H]
\caption{圧力の単位の確認（5data $\to$ 3data。\file{unit_check.csv}）}\label{tab:unit}
\begin{tabularx}{\linewidth}{@{}l>{\raggedright\arraybackslash}p{3\zw}rrrr@{}}
\toprule
材料 & モデル & 予測平均の差の最大 & 予測標準偏差の差の最大 & 切片の差（MPa $-$ Pa） & 係数 $\times\ln 10^6$ \\
\midrule
@@TABLE_UNIT@@
\bottomrule
\end{tabularx}
\end{table}

\subsection{次の手順（3b）}
\begin{points}
\pt{モデル}圧力（なし／\PII）、$\Tff$（なし／あり）、$\Tset$（なし／あり）を組み合わせて比べます。基本の入力は $\ln\gd$ と $\ln H$ で、$\Tff$ が本当に要らないかを、グラフを目で見て確かめます。GPR のカーネルには線形項を足します。
\pt{分け方}5data → 3data（S\_V）、210 ℃ を全部テストにして 190・230 ℃ で学習（S\_T）、ランダム 20\%。
\pt{見るもの}設定温度ごとの予測曲線、190 ℃ と 230 ℃ の予測の比、学習範囲の外での予測の幅。成形条件から見たばらつき（射出速度など事前に決まる値を入力にしたモデル）も調べます。
\end{points}

\end{document}
